%% ══════════════════════════════════════════════════════════════
%%  CHAPTER 14 — PROPOSED IMPROVEMENTS
%% ══════════════════════════════════════════════════════════════
\chapter{Proposed Improvements}
\label{ch:improvements}

\section{AutoScholar v2: Deep Search Architecture}
\label{sec:improve:deepsearch}

\subsection{Motivation}

The current system uses a single-pass approach: one Claude context window handles the entire pipeline.
A ``Deep Search'' architecture would decompose the pipeline into multiple specialized agents, each with their own context window, communicating through a shared state store.

\subsection{Proposed Architecture}

\begin{description}
  \item[Agent Pool.] Maintain a pool of specialized agent instances:
  \begin{itemize}
    \item \textbf{Gauss Agent} (Haiku/Sonnet) --- Lightweight, fast literature search.
    \item \textbf{Turing Agent} (Haiku/Sonnet) --- Parallel data collection across multiple sources.
    \item \textbf{Newton Agent} (Sonnet) --- Data preparation and quality validation.
    \item \textbf{Fisher Agent} (Opus) --- Complex analysis requiring deep reasoning.
    \item \textbf{Euler Agent} (Opus) --- Orchestration and academic writing.
    \item \textbf{Popper Agent} (Sonnet) --- Independent review with no context contamination.
  \end{itemize}

  \item[Shared State Store.] A Redis or PostgreSQL-backed state store holding:
  \begin{itemize}
    \item Discovered papers and citations.
    \item Downloaded datasets and metadata.
    \item Analysis results and figures.
    \item Written sections and BibTeX entries.
  \end{itemize}

  \item[Message Bus.] An event-driven message bus (Redis Streams, NATS) enabling:
  \begin{itemize}
    \item Parallel execution of independent stages.
    \item Dynamic agent scaling based on workload.
    \item Fault isolation between agents.
  \end{itemize}
\end{description}

\subsection{Benefits}

\begin{itemize}
  \item \textbf{Parallel execution:} Gauss and Turing can run concurrently.
  \item \textbf{Cost optimization:} Use cheaper models (Haiku) for simple tasks.
  \item \textbf{Fault isolation:} A Python crash in Fisher doesn't affect Euler.
  \item \textbf{Context efficiency:} Each agent only needs context relevant to its task.
  \item \textbf{Independent review:} Popper reviews without seeing the generation process, reducing self-review bias.
\end{itemize}

\section{Journal++ Model}
\label{sec:improve:journal}

\subsection{Multi-Tier Publication}

Introduce quality tiers for published papers:

\begin{description}
  \item[Tier A (Score $\geq$ 80):] ``Featured Research'' --- Highlighted on the journal homepage, included in monthly digests.
  \item[Tier B (Score 65--79):] ``Published Research'' --- Standard publication, included in domain listings.
  \item[Tier C (Score 45--64):] ``Working Papers'' --- Published with a disclaimer, excluded from quality metrics.
  \item[Rejected (Score $<$ 45):] Archived but not published. Available via admin API for analysis.
\end{description}

\subsection{Multi-Pass Review}

\begin{enumerate}
  \item \textbf{Pass 1:} \popper generates initial review (current system).
  \item \textbf{Pass 2:} A second reviewer (different model or temperature) provides an independent assessment.
  \item \textbf{Pass 3:} Meta-reviewer reconciles the two reviews and makes the final decision.
  \item \textbf{Optional:} Human reviewer for Tier A candidates.
\end{enumerate}

\subsection{Citation Network}

Enable cross-paper citation:
\begin{itemize}
  \item When a new paper is generated, search the existing corpus for relevant prior work.
  \item Auto-generate BibTeX entries for \autoscholar papers.
  \item Track citation relationships in the database.
  \item Display citation graphs on paper landing pages.
\end{itemize}

\subsection{Revision Workflow}

Instead of single-pass publication:
\begin{enumerate}
  \item Generate paper (current pipeline).
  \item Review (Popper).
  \item If score 45--65: trigger a targeted revision cycle.
  \item Euler addresses specific criticisms, re-runs analysis if needed.
  \item Second review.
  \item Publish or reject based on improved score.
\end{enumerate}

\section{Better Memory System}
\label{sec:improve:memory}

\subsection{Cross-Paper Knowledge Base}

Implement a persistent knowledge base that accumulates across paper generations:

\begin{description}
  \item[Topic Memory.] Track topics covered with timestamps. Prevent duplicate topics within a configurable window (e.g., 30 days). Weight prompt selection away from recently covered areas.

  \item[Method Memory.] Track which methods have been applied to which domains. Prefer novel combinations. Detect and avoid repetitive methodological patterns.

  \item[Data Source Memory.] Track which data sources have been used, their reliability, and common issues. Learn which FMP endpoints frequently return errors for specific tickers.

  \item[Failure Memory.] Record common failure patterns with their resolutions. When a similar error occurs, inject the known solution as a governor directive.

  \item[Quality Memory.] Track Popper scores by domain, method, and prompt type. Adjust future prompt selection to favor combinations that historically produce higher scores.
\end{description}

\subsection{Implementation}

\begin{itemize}
  \item Store knowledge base entries in a dedicated database table with vector embeddings for semantic search.
  \item Inject relevant knowledge base entries into the system prompt at the start of each run.
  \item Limit injected context to 5,000 tokens to avoid overwhelming the system prompt.
  \item Use a sliding window to decay old entries and prioritize recent learnings.
\end{itemize}

\section{True Novelty Detection}
\label{sec:improve:novelty}

\subsection{Current Limitation}

The current gap detection relies on Claude's reasoning about a small sample of papers.
True novelty requires systematic comparison against the full academic corpus.

\subsection{Proposed System}

\begin{enumerate}
  \item \textbf{Embedding-based search:} Compute embeddings for all paper abstracts (generated and external) using a dedicated embedding model.
  \item \textbf{Novelty scoring:} Before starting a paper, compute the cosine similarity between the proposed topic and all existing papers. Flag topics with similarity $> 0.9$ to any existing paper.
  \item \textbf{Contribution verification:} After writing, verify that the paper's stated contributions are not already present in the literature by embedding comparison.
  \item \textbf{Plagiarism detection:} Compare generated text against existing papers to detect excessive similarity (beyond proper citation).
\end{enumerate}

\subsection{Technical Requirements}

\begin{itemize}
  \item Vector database (pgvector, Pinecone, or Qdrant).
  \item Embedding model (OpenAI Ada, BGE, or E5 via OpenRouter).
  \item Nightly batch embedding job for new papers.
  \item Similarity threshold tuning based on domain characteristics.
\end{itemize}

\section{Dataset Versioning}
\label{sec:improve:versioning}

\subsection{Problem}

Currently, data is fetched at pipeline runtime, and the raw API responses are not stored.
This means:
\begin{itemize}
  \item Reproducibility requires re-fetching from live APIs.
  \item API changes can break reproduction.
  \item There is no versioning of datasets across papers.
\end{itemize}

\subsection{Proposed System}

\begin{description}
  \item[Data Lake.] Store all fetched API responses in a versioned data lake (S3 or local filesystem with content-addressable storage).

  \item[Dataset Registry.] Maintain a registry of datasets with:
  \begin{itemize}
    \item Content hash (SHA-256) for deduplication.
    \item Source API, parameters, and timestamp.
    \item Data schema (column types, row count).
    \item Usage tracking (which papers reference which datasets).
  \end{itemize}

  \item[Cache Layer.] Before fetching from external APIs, check the cache:
  \begin{itemize}
    \item If data exists and is $<$ 24 hours old, use cached version.
    \item If data exists but is stale, fetch new version and store alongside.
    \item Track all versions for reproducibility.
  \end{itemize}
\end{description}

\section{Reproducibility Layer}
\label{sec:improve:reproducibility}

\subsection{Enhanced Project Archives}

\begin{enumerate}
  \item \textbf{Pin library versions:} Generate \texttt{requirements.txt} with exact version pins (e.g., \texttt{pandas==2.1.4}).
  \item \textbf{Include raw data:} Bundle fetched API responses in the project ZIP.
  \item \textbf{Set random seeds:} Inject \texttt{random.seed(42); np.random.seed(42)} into all scripts.
  \item \textbf{Docker support:} Generate a \texttt{Dockerfile} with the exact Python version and dependencies.
  \item \textbf{GitHub Actions:} Generate a CI pipeline that re-runs the analysis and verifies outputs.
\end{enumerate}

\subsection{Verification Service}

\begin{itemize}
  \item Automated re-execution of project archives in a clean environment.
  \item Comparison of outputs (figures, tables, statistics) against the published paper.
  \item Reproducibility score based on output match percentage.
  \item Public reproducibility badge on paper landing pages.
\end{itemize}

\section{Additional Improvements}
\label{sec:improve:additional}

\subsection{Governor v2: Hybrid Rule-LLM}

Enhance the governor with selective LLM calls for complex decisions:
\begin{itemize}
  \item Use rules for common patterns (current behavior).
  \item Use Claude Haiku for novel error patterns not covered by rules.
  \item Gradually convert successful LLM decisions into new rules.
\end{itemize}

\subsection{Interactive Revision Mode}

Allow human researchers to:
\begin{itemize}
  \item Request specific revisions to generated papers.
  \item Provide additional data or constraints.
  \item Override Popper's review decision.
  \item Add their own analysis as co-authors.
\end{itemize}

\subsection{Multi-Language Support}

\begin{itemize}
  \item Generate papers in French, Spanish, German, Chinese, and Japanese.
  \item Language-specific style guides and citation formats.
  \item Bilingual abstracts.
\end{itemize}

\subsection{Real-Time Data Integration}

\begin{itemize}
  \item WebSocket connections to live market data feeds.
  \item Real-time event detection triggering paper generation.
  \item Breaking news analysis and rapid-response research.
\end{itemize}

\subsection{Performance Optimization}

\begin{itemize}
  \item Python process pooling to amortize startup costs.
  \item Caching layer for frequently accessed API data.
  \item Streaming \LaTeX{} compilation (compile sections as they are written).
  \item Parallel tool execution for independent tool calls.
\end{itemize}
